\chapter{Propozycja testów}
\label{cha:propozycja-testow}

Testowanie systemu OSK-Manager powinno obejmować zarówno logikę biznesową backendu, jak i poprawność działania interfejsu użytkownika. Szczególne znaczenie mają obszary związane z autoryzacją, planowaniem lekcji, dostępnością instruktorów, dostępnością pojazdów, płatnościami oraz ocenami lekcji.

\section{Testy jednostkowe}
\label{sec:testy-jednostkowe}

Testy jednostkowe powinny sprawdzać pojedyncze funkcje i reguły domenowe bez konieczności uruchamiania całej aplikacji. W systemie OSK-Manager szczególnie istotne są testy:
\begin{itemize}
    \item sprawdzania poprawności identyfikatorów UUID,
    \item walidacji statusów pojazdów,
    \item kwalifikacji instruktora do prowadzenia danego typu kursu,
    \item domyślnych godzin pracy instruktora,
    \item zasad rejestracji użytkowników i przypisywania ról,
    \item obsługi błędów integracji z Supabase.
\end{itemize}

Takie testy ograniczają ryzyko przypadkowej zmiany podstawowych reguł systemu podczas dalszego rozwoju aplikacji.

\section{Testy usług backendowych}
\label{sec:testy-uslug}

Najważniejsza logika systemu znajduje się w usługach backendowych. Testy tej warstwy powinny obejmować:
\begin{itemize}
    \item tworzenie i edycję lekcji,
    \item anulowanie lekcji,
    \item rezerwację lekcji przez kursanta,
    \item sprawdzanie konfliktów harmonogramu,
    \item obsługę wydarzeń instruktorskich,
    \item wyliczanie dostępności instruktorów,
    \item obsługę płatności kursanta,
    \item zmianę statusu procesu szkolenia,
    \item wystawianie ocen zakończonych lekcji.
\end{itemize}

Testy powinny weryfikować zarówno scenariusze poprawne, jak i przypadki błędne. Przykładowo system powinien odrzucić próbę zaplanowania lekcji w terminie, w którym instruktor jest zajęty, pojazd jest niedostępny lub dane należą do innego OSK.

\section{Testy API i kontraktu}
\label{sec:testy-api}

Backend udostępnia API HTTP, dlatego należy testować poprawność endpointów oraz zgodność dokumentacji OpenAPI z implementacją. Testy kontraktu API powinny sprawdzać, czy najważniejsze ścieżki są opisane, czy odpowiedzi mają oczekiwaną strukturę oraz czy schematy walidacji odpowiadają rzeczywistym danym zwracanym przez backend.

Warto objąć testami następujące obszary API:
\begin{itemize}
    \item logowanie, odświeżanie sesji i pobieranie danych użytkownika,
    \item operacje na kursantach,
    \item operacje na instruktorach,
    \item operacje na pojazdach,
    \item operacje na kursach i typach kursów,
    \item planowanie i edycję lekcji,
    \item harmonogram szkoły,
    \item oceny lekcji.
\end{itemize}

\section{Testy warstwy BFF}
\label{sec:testy-bff}

Warstwa BFF po stronie Nuxt wymaga osobnych testów, ponieważ odpowiada za pośredniczenie między frontendem i backendem. Testy powinny sprawdzać wybór adaptera komunikacji, poprawność przekazywania żądań do backendu, obsługę błędów połączenia oraz walidację parametrów żądań.

Szczególnie ważne jest sprawdzenie, czy aplikacja poprawnie zachowuje się w przypadku braku konfiguracji backendu, błędnego adresu API lub niedostępnego upstreamu. Użytkownik powinien otrzymać czytelny komunikat błędu zamiast technicznego wyjątku.

\section{Testy manualne GUI}
\label{sec:testy-manualne-gui}

Testy manualne powinny obejmować pełne scenariusze wykonywane z perspektywy użytkownika. Proponowane przypadki testowe to:
\begin{enumerate}
    \item Zalogowanie się do systemu poprawnymi danymi.
    \item Próba logowania z błędnym hasłem.
    \item Utworzenie nowego kursanta przez managera.
    \item Przypisanie kursanta do kursu.
    \item Dodanie instruktora i ustawienie jego dostępności.
    \item Dodanie pojazdu i zmiana jego statusu dostępności.
    \item Zaplanowanie lekcji praktycznej z instruktorem i pojazdem.
    \item Próba zaplanowania lekcji w terminie powodującym konflikt.
    \item Anulowanie lekcji.
    \item Wystawienie oceny zakończonej lekcji przez kursanta.
    \item Sprawdzenie listy płatności kursanta.
    \item Sprawdzenie harmonogramu szkoły w widoku managera.
\end{enumerate}

Każdy test manualny powinien określać warunki początkowe, kroki wykonania oraz oczekiwany rezultat. Dzięki temu testy mogą być powtarzane po każdej większej zmianie w aplikacji.

\section{Testy regresji}
\label{sec:testy-regresji}

Testy regresji powinny być wykonywane po zmianach w logice harmonogramu, autoryzacji, modelu danych oraz formularzach. Szczególnie wrażliwe obszary to planowanie lekcji, sprawdzanie dostępności, przypisywanie danych do OSK oraz uprawnienia użytkowników.

Minimalny zestaw regresji powinien obejmować uruchomienie testów automatycznych backendu i frontendu, sprawdzenie typów TypeScript, lintowanie oraz ręczne przejście najważniejszych scenariuszy GUI. Pozwala to szybko wykryć błędy w miejscach, które nie były bezpośrednio zmieniane, ale zależą od wspólnej logiki systemu.
